% --------------------------------------------------------------------------
\section{Cluster-Level DRO: Discovering Latent Groups}
\label{sec:ch4_attackdropp}

This section introduces the cluster-level DRO component used by AttackDRO++. Instead
of assigning one group to each source attack identity, the component discovers latent
groups from adversarial examples. The central hypothesis is that attack identity is
too coarse as a grouping criterion. Two adversarial examples generated by different
attacks may share the same failure mode, while two examples generated by the same
attack may have very different margins, representations, and optimization directions.
Cluster-based grouping allows the DRO objective to adapt to hard subgroups that may
cut across attack boundaries. The complete AttackDRO++ method is defined after the
augmented clustering features and uniform anchor are introduced in Sections~\ref{sec:ch4_gradfp}
and~\ref{sec:ch4_anchor_objective}.

For each adversarial example \(\mathbf{x}_i^{\mathrm{adv}}\), the model extracts a
representation vector
\begin{equation}
    \mathbf{h}_i = \phi_\theta(\mathbf{x}_i^{\mathrm{adv}}),
\end{equation}
and a clustering model assigns the example to a latent group
\begin{equation}
    c_i \in \{1,\ldots,K\}.
\end{equation}
Let \(\mathcal{B}_k\) denote the subset of the current minibatch whose adversarial
examples are assigned to cluster \(k\):
\begin{equation}
    \mathcal{B}_k
    =
    \{i \in \{1,\ldots,N\}: c_i = k\}.
\end{equation}
The cluster loss is
\begin{equation}
    \mathcal{L}_k
    =
    \frac{1}{|\mathcal{B}_k|}
    \sum_{i\in\mathcal{B}_k}
    \ell\!\left(
        f_\theta(\mathbf{x}_i^{\mathrm{adv}}),
        y_i
    \right).
    \label{eq:cluster_loss}
\end{equation}

The cluster-DRO component maintains a cluster-level weight vector
\begin{equation}
    \mathbf{q} = (q_1,\ldots,q_K) \in \Delta^K.
\end{equation}
Because not every cluster appears in every minibatch, the cluster weights are
renormalized over the clusters present in the current minibatch. Let
\begin{equation}
    \mathcal{C}_{\mathcal{B}}
    =
    \{k \in \{1,\ldots,K\}: |\mathcal{B}_k|>0\}
\end{equation}
denote the set of present clusters. For each \(k \in \mathcal{C}_{\mathcal{B}}\), we
define
\begin{equation}
    \tilde{q}_k
    =
    \frac{q_k}
    {\sum_{j\in\mathcal{C}_{\mathcal{B}}} q_j}.
\end{equation}
The cluster-DRO loss is then
\begin{equation}
    \mathcal{L}_{\mathrm{ClusterDRO}}
    =
    \sum_{k\in\mathcal{C}_{\mathcal{B}}}
    \tilde{q}_k \mathcal{L}_k.
    \label{eq:cluster_dro_loss}
\end{equation}

This objective changes the unit of robustness balancing. Rather than assigning one
weight to each attack identity, the model assigns weight to discovered subgroups
of adversarial examples, so hard regions that cut across attack families can receive
more training emphasis. The renormalization over present clusters keeps the minibatch
objective well-defined even when some latent groups are absent from a particular batch.

The cluster weights are updated using the same exponentiated-gradient principle as in
AttackDRO, but now at the cluster level rather than the attack level. Given a
cluster-level difficulty score
\begin{equation}
    s_k
    =
    \frac{1}{|\mathcal{B}_k|}
    \sum_{i\in\mathcal{B}_k} s_i,
    \qquad k \in \mathcal{C}_{\mathcal{B}},
\end{equation}
the update increases the weight of clusters with higher average difficulty. This allows
the cluster-DRO component to emphasize hard latent groups even when they do not align exactly with
the original attack identities.

Since model representations evolve during training, cluster assignments must also be
updated. The clustering procedure therefore maintains a feature bank \(\mathcal{F}\), which stores
detached augmented feature vectors from recent minibatches. When the refresh condition
is met, K-means is refitted on \(\mathcal{F}\), and subsequent minibatches are assigned
using the updated clusterer. Before the first K-means refresh, cluster assignments are
bootstrapped randomly.

The refresh condition is treated as a policy rather than a fixed part of the method. It
can be defined in terms of epochs or training steps. In the main experiments, we use an
epoch-based refresh schedule with \(R=2\) epochs. After each refresh, the feature bank
is reset or updated according to the memory policy, and the cluster weights are refreshed
according to the chosen \(q\)-refresh policy.
